\itemhead{2.1 例题 · 例题9（教材p.41）}{XJ-B1-C02-EX-P041-Y1378}
\label{XJ-B1-C02-EX-P041-Y1378}
{\footnotesize 湘教版教材p.41 / PDF第48页 / 2.1 例题 / 例题9}
\field{原题}{某公司设计了如图 2.1-5 所示的一块绿化景观地带，两条平行线段的两端用半圆弧相连接。已知这块绿化景观地带的内圈周长为 $400\,\mathrm m$，当平行线段的长设计为多少时，中间矩形区域的面积最大？
\bankfig{C02-EX041}{0.46}}
\field{参考答案}{设平行线段长为 $l$、两线间距为 $w$，则 $2l+\pi w=400$。$lw=l(400-2l)/\pi$ 在 $l=100$ 米时最大，为 $20000/\pi$ 平方米。}
{\small \field{关联知识}{I03 基本不等式、等号与最值建模（教材帮p.52、53、54 / PDF61、62、63）；F03 函数求值、解析式与函数建模（教材帮p.80、87、88、89、90、131 / PDF89、96、97、98、99、140）}
\field{方法与易错点}{列正数条件与固定量，构造基本不等式，解等号条件并代回实际限制。 易错：把一个下界误当最小值；多个等号条件不能同时实现；漏单位和整数限制。}
\field{答案依据}{例题原解位置：湘教版教材p.41及题库原图；独立编写可编辑参考解答。答案册此部分主要收练习、习题，不伪造例题答案册位置。}}
